\section{Introduction}
\label{sec:intro}

% 证据：调研笔记 01、03
% 段落计划：
%  1. 间接编码 + 解码器在调度/排班中很常见；解码器带来多对一（冗余）。
%  2. 规则（启发式）偏置有用但不稳健——好解离规则远时偏置有害（Rothlauf 的冗余/偏置理论）。
%  3. 本文问题：如何同时给出偏置与自由度，以及如何获得"好的"规则。
%  4. 本文做法概述。
%  5. 贡献列表（见下）。
%  6. 文章结构。
% 图 1：编码示意（dec5 逐班基因 vs 规则 + 稀疏偏离）

\todo{Introduction.}

The main contributions of this paper are as follows.
\begin{enumerate}
  \item \textbf{Representation.} \todo{Rules plus sparse deviation slots, with the deviation density as an explicit parameter (LDS idea turned into an MOEA encoding).}
  \item \textbf{Theory.} \todo{Bias--freedom model: for $\K<\dist$ search stalls above $\dist-\K$; small-budget optimum $\K\approx\dist$, large-budget $\approx\dist/\nalt$; multi-objective case follows the mean distance.}
  \item \textbf{Rule acquisition.} \todo{Rules should be optimized for bias quality (the proxy score); under this objective GP, neural network + ES and LLM are interchangeable; imitating good solutions is not enough.}
  \item \textbf{Applications.} \todo{Static rostering (47 test instances), mid-execution re-rostering, and public FJSP benchmarks.}
\end{enumerate}

% Organization paragraph
\todo{Organization of the paper.}
